\documentclass[12pt]{report}
\usepackage{xcolor}
\usepackage[utf8]{inputenc}
\usepackage{natbib}%referencing
\usepackage[left=1.0in,right=1.0in,top=.8in,bottom=.8in]{geometry}
\usepackage{float}
\linespread{1.3}
\usepackage{graphicx}%figures
\usepackage{rotating}%landscape
\usepackage{amsmath}%math
\usepackage{titlesec} %formatting chapters
\usepackage{svg}
%\titlespacing*{<command>}{<left>}{<before-sep>}{<after-sep>}
\titlespacing*{\chapter}{-15pt}{10pt}{15pt}
\titlespacing*{\section}{0pt}{0pt}{5pt}
\titlespacing*{\subsection}{0pt}{5pt}{5pt}
\titleformat{\chapter}[hang] 
{\normalfont\huge\bfseries}{\chaptertitlename\ \thechapter.}{1em}{}
\renewcommand{\chaptername}{}
\graphicspath{{Images/}}%image folder name
%Cover page contents
\title{
{\includegraphics[scale=.3]{logotu.jpg}}\\
\uppercase\large{
    {Tribhuvan University}\\
    {Institute of Engineering}\\
    {Pulchowk Campus}\\
    \vspace{.5cm}
    {A \\Project Report\\On\\Title of project}\\
    \vspace{.5cm}
    {\textbf{Submitted By:}\\Niraj Karki (PUL079BCT049) \\Pariskar Sharma Poudel (PUL079BCT055) \\Suyash Adhikari (PUL079BCT090)\\
    \vspace{.5cm}
    {\textbf{Submitted To:}\\ Department of Electronics \& Computer engineering}\\
                }
    }}
\datetime %{Month Year}

\begin{document}
\maketitle
\pagenumbering{roman}
\setcounter{page}{2}
\chapter*{Page of Approval}
\addcontentsline{toc}{chapter}{\numberline{}Page of Approval}
\vspace{1.5cm}
\begin{center}
\uppercase{
Tribhuvan Universiy\\
Institute of Engineering\\
Pulchowk Campus\\
Department of Electronics and Computer Engineering
    }
    \end{center}
\vspace{.5cm}
The undersigned certifies that they have read and recommended to the Institute of Engineering for acceptance of a project report entitled \textbf{``PINN: ''} submitted by \textbf{Niraj Karki}, \textbf{Suyash Adhikari}, \textbf{Pariskar Sharma Poudel} in partial fulfillment of the requirements for the Bachelor's degree in Electronics \& Computer Engineering.\\

\vspace{1cm}
\begin{minipage}{.5\textwidth}
        \raggedright
    .............................\\
    Supervisor\\
    \textbf{Sanjeeb Prasad Pandey}\\
    Designation\\
    Department of Electronics and Computer Engineering,\\
    Pulchowk Campus, IOE, TU.\\
    \end{minipage}%
    \begin{minipage}{0.5\textwidth}
    \raggedleft
    .............................\\
    External examiner\\

    Assistant Professor\\
    Department of Electronics and Computer Engineering,\\
    Pulchowk Campus, IOE, TU.\\
    \end{minipage}\\
 \newline
 \newline
 \newline
\centerline{Date of approval:}
%new chapter
\chapter*{Copyright}
\addcontentsline{toc}{chapter}{\numberline{}Copyright}
The author has agreed that the Library, Department of Electronics and Computer Engineering, Pulchowk Campus, and Institute of Engineering may make this report freely available for inspection. Moreover, the author has agreed that permission for extensive copying of this project report for scholarly purposes may be granted by the supervisors who supervised the work recorded herein or, in their absence, by the Head of the Department wherein the project report was done. It is understood that recognition will be given to the author of this report and the Department of Electronics and Computer Engineering, Pulchowk Campus, Institute of Engineering for any use of the material of this project report. Copying publication or the other use of this report for financial gain without the approval of to the Department of Electronics and Computer Engineering, Pulchowk Campus, Institute of Engineering, and the author's written permission is prohibited.\\
Request for permission to copy or to make any other use of the material in this report in whole or in part should be addressed to:\\
\newline
\newline
\newline
Head\\
Department of Electronics and Computer Engineering\\
Pulchowk Campus, Institute of Engineering, TU\\
Lalitpur, Nepal.
%new chapter
\chapter*{Acknowledgments}
\addcontentsline{toc}{chapter}{\numberline{}Acknowledgements}
First and foremost, we would like to express our sincere indebtedness to the Department of Electronics and Computer Engineering at Pulchowk Campus for granting us the opportunity to undertake this research. The department's consistent support, meticulous administrative assistance, and commitment to fostering an environment of academic excellence have been integral to the successful execution of this study.

We wish to extend our utmost gratitude and deepest respect to our supervisor, Prof. Dr. Sanjeeb Prasad Pandey, whose exceptional guidance, profound expertise, and unwavering encouragement have played an instrumental role in shaping the trajectory and outcomes of this research endeavor. His insightful direction during our regular progress meetings was invaluable in navigating the technical complexities of Physics-Informed Neural Networks (PINNs) and their application to landslide modeling. We are truly honored to have worked under his distinguished mentorship.

%new chapter
\chapter*{Abstract}

Rainfall-induced landslides claim hundreds of lives annually across Nepal's Himalayan region. Because real-time tracking of subsurface pore-water pressure is operationally infeasible for most unmonitored slopes, this project develops a Hybrid Physics-Informed Neural Network (PINN) to act as a virtual sensor. Trained on physically consistent synthetic data, the model evaluates slope stability at any depth and time in milliseconds without requiring in-situ field instrumentation.

The framework is applied to Nepal's catastrophic 2014 Jure landslide. HYDRUS-1D serves as a teacher model, simulating 1D vertical unsaturated infiltration governed by the Richards equation to generate 24,000 pseudo-data points of pressure head $\psi(z,t)$. The PINN is trained using a physics-aware loss function that enforces the Richards equation via automatic differentiation, satisfies boundary conditions, and aligns with anchor data using Adam and L-BFGS optimisation. The Factor of Safety is computed using the infinite-slope stability model with Van Genuchten--Mualem hydraulic properties.

The trained PINN achieved a 99.99\% reduction in the Richards equation residual and reproduced the reference pressure head field with a relative RMSE of 0.16\%. It successfully identified the critical stability transition depth ($z = 1.16$--$1.19$ m) within 4.1\% of the HYDRUS reference. Evaluating 2,880 query points in under 50 milliseconds—three to four orders of magnitude faster than traditional solvers—the model is deployed via an interactive dashboard and REST API. This framework establishes a technical foundation for low-cost, real-time early warning systems and transfer learning across the unmonitored Hindu Kush Himalayan region.

\vspace{0.5cm}
\noindent\textbf{Keywords:} \textit{Physics-Informed Neural Networks, Richards equation, rainfall-induced landslides, slope stability, virtual sensor}

\addcontentsline{toc}{chapter}{\numberline{}Abstract}
\tableofcontents
\addcontentsline{toc}{chapter}{\numberline{}Contents}
\listoffigures
\addcontentsline{toc}{chapter}{\numberline{}List of Figures}
\listoftables
\addcontentsline{toc}{chapter}{\numberline{}List of Tables}
\chapter*{List of Abbreviations}
\addcontentsline{toc}{chapter}{\numberline{}List of Abbreviations}
\begin{tabular}{c l}
\textbf{LAH}     &  List of Abbreviations Here\\
\textbf{AMC}     &  Add More Contents
\end{tabular}
\chapter{Introduction}
\pagenumbering{arabic}%start arabic numbering(1,2,3..) from here
\section{Background}
{ Nepal is characterized by its rugged Himalayan topography and intense seasonal weather patterns. During the South Asian Monsoon (June to September), the country receives approximately 80\% of its total annual precipitation. Consequently, Nepal faces a chronic threat from rainfall-induced landslides, which result in a devastating annual toll on human life and national infrastructure (Dahal and Hasegawa 2008).}
\newline
\indent {
The initiation of these landslides is primarily driven by rainwater infiltration into the unsaturated soil mantle. As soil moisture increases, the accompanying rise in pore-water pressure reduces the effective shear strength of the slope, eventually leading to structural collapse. To predict such failures, it is essential to monitor the dynamic hydraulic state, specifically the pressure head and moisture content - within the subsurface.}
\newline
\indent {
Traditionally, tracking these processes requires expensive geotechnical instrumentation, such as tensiometers or piezometers. However, for a developing nation like Nepal, the installation and long-term maintenance of dense sensor networks across thousands of hazardous, remote slopes is logistically and economically unfeasible. While Artificial Intelligence (AI) has emerged as a powerful alternative for disaster forecasting, standard "Black-Box" deep learning models are notoriously data-hungry. They require massive historical datasets that simply do not exist for the vast majority of unmonitored slopes in the Himalayas.}
\newline
\indent{
To bridge this gap, this research proposes a hybrid approach using Physics-Informed Neural Networks (PINNs). By embedding fundamental physical laws—specifically the Richards Equation for unsaturated flow—directly into the neural network's loss function, PINNs can simulate complex subsurface dynamics using only surface-level rainfall data. This framework allows the model to ``learn'' from limited data while strictly adhering to the laws of physics, effectively functioning as a ``virtual sensor'' for data-scarce regions (Raissi et al. 2019).
}


\section{Problem Statement}

The fundamental engineering challenge in landslide risk reduction in
Nepal is the critical trade-off between data scarcity and physical
accuracy. A reliable Early Warning System (EWS) necessitates real-time
tracking of pore-water pressure evolution within a slope; however,
achieving this in the Himalayan context is obstructed by two primary
barriers:

\begin{itemize}

    \item \textbf{The Data Barrier:} In-situ subsurface monitoring
    records, such as soil matric suction and volumetric water content,
    are virtually non-existent for most hazardous slopes in Nepal.
    Without these high-resolution historical records, conventional
    ``Black-Box'' Artificial Intelligence models lack the ground truth
    required to map rainfall patterns to slope instability accurately.

    \item \textbf{The Computational Barrier:} Traditional numerical
    frameworks such as HYDRUS-1D and Finite Element Analysis can
    simulate subsurface flow using physical parameters alone, but are
    computationally intensive and unsuitable for rapid real-time
    deployment — a single HYDRUS scenario requires 30--120 seconds,
    making continuous slope monitoring operationally infeasible.

\end{itemize}

This project addresses both barriers simultaneously by developing a
hybrid PINN solver that functions as a \textbf{virtual sensor} for
unmonitored slopes — trained once on HYDRUS-generated pseudo-data to
embed the governing physics, then deployed to evaluate pore-water
pressure and slope stability at any depth and time in under 50
milliseconds, without requiring any in-situ instrumentation.
\section{Objectives}

The primary objective of this project is to develop a Hybrid
Physics-Informed Neural Network (PINN) Framework that solves the
Richards Equation to simulate unsaturated soil moisture dynamics and
characterise slope stability in data-scarce regions, using the 2014
Jure landslide as the primary case study. The specific objectives are:

\begin{itemize}

    \item To implement a PINN architecture that integrates the
    \textbf{Richards Equation} and the \textbf{Van Genuchten--Mualem}
    model as governing physical constraints enforced through automatic
    differentiation.

    \item To develop a hybrid training strategy using HYDRUS-1D
    numerical pseudo-data as anchor points to calibrate the model for
    site-specific hydraulic conditions of the Jure slope.

    \item To predict the spatiotemporal pressure head field
    $\psi(z,t)$ and compute the Factor of Safety using the
    infinite-slope stability model, characterising the critical
    stability transition depth of the slope.

    \item To demonstrate the framework as a production-ready virtual
    sensor through an interactive dashboard providing real-time
    inference, sensitivity analysis, and uncertainty quantification.

\end{itemize}

\section{Scope}

This project implements a complete end-to-end Hybrid PINN pipeline for
rainfall-induced landslide analysis applied to the 2014 Jure landslide,
Sindhupalchok, Nepal. The framework models 1D vertical unsaturated
infiltration over a 55 m soil column using HYDRUS-1D-generated
pseudo-data (24,000 points across 123 days) as anchor data, directly
addressing the data-scarcity challenge of unmonitored Himalayan slopes.

The trained PINN approximates the continuous pressure head field
$\psi(z,t)$ using a 7-layer, 64-neuron-per-layer architecture with
a five-component physics-aware loss function, and computes the Factor
of Safety at any queried depth and time in under 50 milliseconds — three
to four orders of magnitude faster than the HYDRUS solver it was trained
on. Results are delivered through a 17-page interactive dashboard with
real-time parameter exploration and Monte Carlo uncertainty
quantification.

The scope is bounded to 1D vertical infiltration physics. Lateral flow,
the inverse parameter identification problem, and extension to other
Himalayan slopes are identified as direct future work enabled by the
framework established here.


\chapter{Literature Review}This chapter reviews the theoretical foundations and recent advances in Physics-Informed Neural Networks (PINNs) for landslide prediction. It explores the transition from traditional geotechnical modeling to hybrid deep learning, highlighting how PINNs integrate physical laws with data-driven learning to improve prediction accuracy in data-scarce regions.\section{Traditional and Data-Driven Methods}\subsection{Limit Equilibrium and Seismic Analysis}Limit Equilibrium Methods (LEM) evaluate slope stability by balancing driving and resisting forces. The Bishop Simplified Method \citep{bishop1955use} remains a standard for calculating the factor of safety ($F_s$):\begin{equation}F_s = \frac{\sum \frac{1}{m_{\alpha_i}} \left[ c_i b_i + (W_i - u_i b_i) \tan \phi_i \right]}{\sum W_i \sin \alpha_i}\end{equation}
\newline
\indent{
While LEM is practical, it assumes simplified failure geometries and provides static estimates \citep{michalowski2009three}. For dynamic triggers, Newmark’s sliding block model \citep{newmark1965effects} estimates seismic displacement by treating the landslide mass as a rigid block. However, these traditional approaches cannot easily model the progressive evolution of pore-water pressure without dense instrumentation.}
\subsection{Machine Learning Approaches}Machine learning methods, such as Random Forest and SVM, have achieved strong performance in regional susceptibility mapping \citep{chowdhury2024landslide}. Deep learning models like CNNs further improve nonlinear pattern learning using satellite datasets. However, their "Black-Box" nature often leads to results that lack physical consistency, limiting their adoption in safety-critical geotechnical applications where interpretability is paramount \citep{von2021black}.\section{Physics-Informed Neural Networks for Landslides}PINNs integrate governing physical laws directly into the neural network training process through automatic differentiation \citep{raissi2019physics}. This paradigm addresses the limitations of purely empirical models by enforcing physical constraints during learning.\subsection{Coseismic and Rainfall-Induced Modeling}\citet{dahal2024towards} applied PINNs to coseismic landslide prediction by embedding Newmark-based physical constraints. Their framework retrieves mechanical parameters from proxy variables, producing both susceptibility maps and spatial distributions of soil properties.For hydrologically triggered failures, PINNs embed the Richards Equation into the loss function to simulate unsaturated infiltration:\begin{equation}\frac{\partial \theta}{\partial t} = \nabla \cdot \left[ K(\theta) \nabla (\psi + z) \right]\end{equation}where $\theta$ denotes volumetric water content and $\psi$ represents pressure head. This enables PINNs to function as surrogate models for hillslope hydrology, allowing for the simulation of pore-pressure dynamics even under sparse observations \citep{depina2022physics, krkač2023physics}.\subsection{Hybrid and Inverse Modeling}To address computational costs, hybrid physics–ML approaches have been explored. Comparative studies indicate that while PINNs may be slower than traditional solvers for simple forward problems, they are highly effective for solving Inverse Problems—discovering unknown hydraulic parameters from past failure data \citep{wang2025physics, niu2025inverse}. Using numerical baseline data as "anchor points" further stabilizes training in data-scarce environments.\section{Research Gaps and Relevance to Nepal}Despite significant advances, gaps remain in computational efficiency and the application of PINNs to regional-scale heterogeneity \citep{wang2025physics}. In Nepal, the Himalayan terrain presents unique challenges: extreme topographic relief, complex monsoon patterns, and a lack of subsurface monitoring infrastructure \citep{adhikari2025ai}.The application of PINNs in Nepal offers a specific solution to Data Scarcity. By using physical laws to "fill the gaps" in missing sensor data, PINNs can infer mechanical and hydraulic properties from rainfall and topographic proxies. This research seeks to bridge the identified gaps by applying a hybrid PINN framework to the 2014 Jure landslide, providing a proof-of-concept for a low-cost Early Warning System in the region.

\setlength{\parindent}{0pt}

\chapter{Study Area}

The Jure landslide (27.46°N, 85.52°E) struck on August 2, 2014, around 2:30 AM in Mankha VDC, Sindhupalchok District, Nepal. The site lies along the right bank of the Sunkoshi River near Jure village, approximately 80 km northeast of Kathmandu along the Araniko Highway. This catastrophic event killed 156 people, destroyed over 120 houses. The landslide mobilized an estimated 5--6 million m³ of debris that traveled approximately 1.3 km downslope, completely damming the Sunkoshi River. The resulting lake ($\sim$3 km long, 300--350 m wide, $\sim$8 million m³ volume) posed significant downstream flood risk until controlled drainage operations were completed.

\section{Geological and Topographic Setting}

The landslide is located in the Lesser Himalaya within the Kuncha Formation (Mesoproterozoic). The lithology is dominated by weathered phyllite, chloritic schist, and quartzite. These low-grade metamorphic rocks are characterized by high foliation and significant fracturing, with four primary joint sets identified in the scar area. The pre-failure slope was exceptionally steep, with angles ranging between 40° to 60°. The presence of near-vertical tension cracks and joints is a critical feature of this site, as they served as preferential pathways for rapid water infiltration of rainfall to depth. The estimated failure surface lies approximately 20--30 m below the ground surface, indicating a deep-seated failure mechanism.

\section{Hydrological Context}

The region experiences a monsoon-dominated climate, with 80\% of the annual precipitation (approx. 3,500 mm) falling between June and September. The 2014 failure was preceded by 60 days of continuous antecedent rainfall.
\begin{itemize}
    \item \textbf{Antecedent Phase:} Cumulative rainfall from mid-June through July exceeded 1,000 mm, causing progressive saturation of the rock mass.
    
    \item \textbf{Trigger Event:} High-intensity bursts of over 140 mm on July 30--31 acted as the immediate trigger, spiking pore-water pressure and reducing the effective shear strength of the weathered phyllite to a critical state (FS $\leq$ 1).

\end{itemize}

\section{Site Selection Rationale}

The Jure landslide is an ideal case study for a Physics-Informed Neural Network (PINN) due to following reasons:

\begin{itemize}
    \item It is a deep-seated failure where 1D vertical infiltration physics (Richards Equation) can accurately model the pore-pressure buildup at depth.
    
    \item It represents a data-scarce environment typical of the Himalayas, where real-time sensor data (suction/moisture) was unavailable, motivating the use of a ``Virtual Sensor'' Hybrid PINN approach.
\end{itemize}

\chapter{Data Acquisition and Sources}

The methodology is designed to function within the constraints of
data-scarce mountainous regions. It relies on a hybrid data strategy
that fuses satellite-derived meteorological records with
literature-derived geotechnical parameters without the need for
in-situ instrumentation.

\section{Dynamic Hydrological Data (Rainfall Time Series)}

\begin{itemize}

    \item \textbf{Source:} PERSIANN-CDR (Precipitation Estimation from
    Remotely Sensed Information using Artificial Neural
    Networks -- Climate Data Record), accessed via the CHRS Data Portal,
    Center for Hydrometeorology and Remote Sensing (CHRS), University
    of California, Irvine
    (\texttt{https://chrsdata.eng.uci.edu}).

    \item \textbf{Dataset:} PERSIANN-CDR provides bias-adjusted daily
    precipitation estimates derived from geostationary satellite
    infrared data, available from 1983 to present at a spatial
    resolution of 0.25° and a temporal resolution of daily totals
    (Ashouri et al., 2015).

    \item \textbf{Extraction point:} Gridded precipitation data
    extracted at the Jure landslide coordinates (27.46°N, 85.52°E),
    Sindhupalchok District, Nepal.

    \item \textbf{Temporal window:} June 1, 2014 -- August 31, 2014
    (full 2014 monsoon season), capturing the 60-day antecedent
    saturation buildup that progressively primed the slope for failure
    prior to the August 2 event.


\end{itemize}
\section{Static Geotechnical and Hydraulic Parameters}

Static parameters for the soil-rock mass are extracted from specialized back-analyses of the Jure event and regional studies of the Kuncha Formation. The key parameters required for the Hydrus-1D numerical simulation (via the Richards equation) and subsequent PINN training are as follows:

\begin{itemize}
    \item Saturated hydraulic conductivity ($K_s$)
    \item Saturated water content ($\theta_s$)
    \item Residual water content ($\theta_r$)
    \item Van Genuchten alpha ($\alpha$)
    \item Pore-size Distribution Index ($n$)
    \item Tortuosity / Connectivity ($l$)
    \item Effective cohesion ($c'$)
    \item Effective friction angle ($\phi'$)
    \item Bulk unit weight ($\gamma$)
\end{itemize}
\newline
\indent{
Primary sources include Bray et al. (2023), Tien et al. (2021), and Panthi (2021). Where site-specific measurements are unavailable, representative regional values for phyllite--schist sequences of the Lesser Himalaya are adopted. To account for uncertainty, a sensitivity analysis is performed by varying selected parameters within $\pm$20--30\% of their nominal values.
}

\section{Topographic and Geometric Characterisation}

To construct the 1D vertical domain the following spatial data is
utilised:

\begin{itemize}
    \item \textbf{Slope angle ($\beta$):} 42° (average in the critical
    initiation zone; range 40°--60° across the failure scar).

    \item \textbf{Profile depth ($z$):} 55 m vertical column
    ($z = 0$ at surface, positive downward), discretised into 1,000
    depth nodes at 0.055 m spacing to resolve sharp near-surface
    moisture gradients during rainfall infiltration.
\end{itemize}

This 1D domain ($z \in [0, 55]$ m) serves as the spatial coordinate
for both HYDRUS-1D and PINN inputs.

\begin{table}[htbp]
\centering
\begin{tabular}{llll}
\hline
\textbf{Category} & \textbf{Parameter} & \textbf{Symbol} &
\textbf{Adopted Value} \\
\hline
\multirow{6}{*}{\textbf{Hydraulic}}
    & Saturated Hydraulic Conductivity & $K_s$       & $1.5\times10^{-5}$ m/s \\
    & Saturated Water Content          & $\theta_s$  & 0.45 \\
    & Residual Water Content           & $\theta_r$  & 0.05 \\
    & Van Genuchten Alpha              & $\alpha$    & 0.8 m$^{-1}$ \\
    & Pore-size Distribution Index     & $n$         & 1.4 \\
    & Tortuosity / Connectivity        & $l$         & 0.5 \\
\hline
\multirow{3}{*}{\textbf{Geotechnical}}
    & Effective Cohesion               & $c'$        & 5.0 kPa \\
    & Effective Friction Angle         & $\phi'$     & 28° \\
    & Bulk Unit Weight                 & $\gamma$    & 20.0 kN/m$^3$ \\
\hline
\multirow{3}{*}{\textbf{Topographic}}
    & Slope Angle                      & $\beta$     & 42° \\
    & Vertical Depth                   & $z$         & 0--55 m \\
    & Unit Weight of Water             & $\gamma_w$  & 9.81 kN/m$^3$ \\
\hline
\textbf{Initial Conditions}
    & Initial Pressure Head            & $\psi(z,0)$ & $-$500 to $-$50 cm \\
\hline
\end{tabular}
\caption{Summary of model parameters for hydraulic and stability
analysis with adopted values for the Jure case study.}
\label{tab:model_parameters}
\end{table}
\chapter{Methodology}

\section{Overview of the Approach}

This study implements a hybrid Physics-Informed Neural Network (PINN)
framework to simulate rainfall-induced subsurface pressure head dynamics
and characterise slope stability in data-scarce Himalayan environments,
using the 2014 Jure landslide as the primary case study. The method
integrates numerical simulation with physics-constrained deep learning
to overcome the absence of real-time field measurements.

The workflow is structured in four phases:

\begin{enumerate}
    \item \textbf{Numerical Baseline:} HYDRUS-1D simulates 1D vertical
    unsaturated flow over the full 2014 monsoon period, generating
    24,000 physically consistent pseudo-data points of pressure head
    $\psi(z,t)$ to anchor the PINN.

    \item \textbf{Data Preparation:} The HYDRUS output is parsed,
    normalised, and assembled into a training batch comprising anchor
    points, collocation points sampled via Latin Hypercube Sampling,
    initial condition points, surface boundary points, and failure
    points.

    \item \textbf{PINN Training:} A 7-layer feed-forward network
    approximates $\psi(z,t)$ by minimising a five-component
    physics-aware loss function. Training uses a two-phase Adam
    + L-BFGS optimisation strategy.

    \item \textbf{Inference and Deployment:} The trained model acts
    as a virtual sensor, evaluating $\psi(z,t)$ and the Factor of
    Safety at any queried depth and time in under 50 milliseconds,
    delivered through a 17-page interactive dashboard.
\end{enumerate}

\section{Hydrogeological Numerical Simulation (HYDRUS-1D)}

A high-fidelity numerical baseline is established using HYDRUS-1D, a
finite-difference solver for the 1D Richards equation governing
unsaturated flow. HYDRUS-1D serves as the ``teacher model'', generating
physically consistent synthetic profiles of pressure head $\psi(z,t)$
that anchor the subsequent PINN training in the absence of any in-situ
sensor measurements.

\subsection{Model Setup}

\textbf{Hydraulic parameters:} The soil water-retention characteristics
are defined using the Van Genuchten--Mualem model with the parameters
listed in Table~\ref{tab:model_parameters}, sourced from back-analyses
of the Jure event (Bray et al., 2023; Tien et al., 2021; Panthi, 2021).

\textbf{Initial state:} The slope is initialised in a dry pre-monsoon
condition with a linear pressure head profile of $\psi \approx -500$
cm at the surface, reflecting typical Himalayan mid-hill dryness prior
to monsoon onset.

\textbf{Boundary conditions:}
\begin{itemize}
    \item \textbf{Top (Atmospheric):} Daily precipitation from the
    PERSIANN-CDR dataset drives the surface infiltration flux $I(t)$.

    \item \textbf{Bottom (Free Drainage):} A unit-gradient condition
    is applied at the base of the 55 m column, simulating
    gravity-driven outflow into deeper permeable bedrock.
\end{itemize}

\subsection{Domain Discretisation and Data Extraction}

The simulation models a 1D vertical column of 55 m ($z = 0$ at
surface, positive downward), discretised into 1,000 depth nodes at
0.055 m uniform spacing. The Richards equation is solved over the full
2014 monsoon period (June 1 -- August 31, 2014), producing 24
daily timesteps.

From the converged simulation, 24,000 pseudo-data points are extracted
— one value of $\psi(z_i, t_j)$ per depth node per timestep. These
points constitute the anchor dataset used to constrain the PINN. The
Factor of Safety is also computed at every point using the
infinite-slope model (Section~\ref{sec:fs}), identifying 23,448 points
where FS $\leq$ 1.0 — used as the failure training set.

\section{Physics-Informed Neural Network Architecture}

The PINN approximates the continuous pressure head field $\psi(z,t)$
as a function of normalised depth and time coordinates. The
implementation uses PyTorch with full automatic differentiation.

\subsection{Network Structure}

The network is a fully connected feed-forward architecture with the
following specification:

\begin{itemize}
    \item \textbf{Input layer:} 2 neurons — normalised depth
    $\tilde{z}$ and normalised time $\tilde{t}$

    \item \textbf{Hidden layers:} 7 fully connected layers, each
    containing 64 neurons with hyperbolic tangent (tanh) activation

    \item \textbf{Output layer:} 1 neuron — normalised pressure head
    $\tilde{\psi}$

    \item \textbf{Weight initialisation:} Xavier normal for weights,
    zeros for biases — standard for PINNs with tanh activation

    \item \textbf{Total parameters:} 29,377
\end{itemize}

\subsection{Input Normalisation}

All inputs and outputs are normalised to $[0, 1]$ before entering the
network to ensure numerical stability and prevent vanishing gradients:

\begin{equation}
\tilde{z} = \frac{z}{z_{\max}}, \quad
\tilde{t} = \frac{t}{t_{\max}}, \quad
\tilde{\psi} = \frac{\psi - \psi_{\min}}{\psi_{\max} - \psi_{\min}}
\end{equation}

where $z_{\max} = 55.0$ m, $t_{\max} = 123.0$ days,
$\psi_{\min} = -500.0$ m, and $\psi_{\max} = -46.91$ m
(range = 453.09 m), derived directly from the HYDRUS training data.
De-normalisation is applied inside the loss function to recover
physical units before computing PDE residuals and hydraulic properties.

\subsection{Collocation and Anchor Points}

The network is trained using four distinct sets of input points:

\begin{itemize}
    \item \textbf{Collocation points ($N_c = 10{,}000$):} Unlabelled
    points generated via Latin Hypercube Sampling (LHS) across the
    normalised domain $[0,1]^2$, used to enforce the Richards equation
    residual. LHS with seed 42 ensures uniform stratified coverage of
    the entire $(z,t)$ domain, avoiding clustering that occurs with
    uniform random sampling.

    \item \textbf{Anchor points ($N_a = 24{,}000$):} All HYDRUS-1D
    pseudo-data points carrying known target values of $\tilde{\psi}$,
    used for the data-driven loss term.

    \item \textbf{Initial condition points ($N_i = 1{,}000$):}
    All depth nodes at $t = t_{\min}$, enforcing the pre-monsoon
    pressure head profile $\psi(z, 0) = \psi_{\text{HYDRUS}}(z, 0)$.

    \item \textbf{Boundary points ($N_b = 24$):} All timesteps at
    $z = z_{\min}$, enforcing the surface flux boundary condition.

    \item \textbf{Failure points ($N_f = 23{,}448$):} All HYDRUS
    records where FS $\leq 1.0$ and depth $> 0.1$ m, used for the
    event-based failure loss.
\end{itemize}

\section{Multi-Component Loss Function}

The PINN minimises a weighted five-component loss:

\begin{equation}
\mathcal{L}_{\text{total}} =
  \lambda_{\text{physics}}  \mathcal{L}_{\text{physics}}  +
  \lambda_{\text{anchor}}   \mathcal{L}_{\text{anchor}}   +
  \lambda_{\text{initial}}  \mathcal{L}_{\text{initial}}  +
  \lambda_{\text{boundary}} \mathcal{L}_{\text{boundary}} +
  \lambda_{\text{failure}}  \mathcal{L}_{\text{failure}}
\end{equation}

Table~\ref{tab:loss_weights} lists the adopted weights and the
justification for each.

\begin{table}[htbp]
\centering
\begin{tabular}{lrp{7cm}}
\hline
\textbf{Component} & \textbf{$\lambda$} & \textbf{Justification} \\
\hline
Physics  & $1 \times 10^6$ & Raw Richards residual is O($10^{-13}$)
                              for dry unsaturated soil; weight
                              compensates after residual normalisation \\
Anchor   & 10.0 & Primary ground-truth signal; highest data weight \\
Initial  & 10.0 & IC as important as anchor; same data type \\
Boundary & 1.0  & Surface flux uses placeholder $-K_s$; low weight
                  until real rainfall record is incorporated \\
Failure  & 20.0 & Primary geomechanical validation target \\
\hline
\end{tabular}
\caption{Loss component weights and justifications.}
\label{tab:loss_weights}
\end{table}

\subsection{Physics Loss ($\mathcal{L}_{\text{physics}}$)}

The physics loss enforces the 1D Richards equation at every
collocation point. The network predicts $\tilde{\psi}$ at each point;
automatic differentiation computes
$\partial\tilde{\psi}/\partial\tilde{z}$,
$\partial\tilde{\psi}/\partial\tilde{t}$, and
$\partial^2\tilde{\psi}/\partial\tilde{z}^2$,
which are de-normalised to physical units before substitution:

\begin{equation}
C(\psi)\frac{\partial \psi}{\partial t} =
\frac{\partial}{\partial z}\left[K(\psi)
\left(\frac{\partial \psi}{\partial z} + 1\right)\right]
\end{equation}

where $C(\psi) = d\theta/d\psi$ is the specific moisture capacity and
$K(\psi)$ is the unsaturated hydraulic conductivity, both evaluated
using the Van Genuchten--Mualem model. The residual is normalised by
a characteristic scale $\xi = C_{\text{char}} \cdot \Delta\psi /
t_{\max,s} = 1.18 \times 10^{-8}$ s$^{-1}$ before squaring, bringing
it to $O(1)$ and making $\lambda_{\text{physics}} = 10^6$ effective:

\begin{equation}
\mathcal{L}_{\text{physics}} = \frac{1}{N_c}
\sum_{j=1}^{N_c} \left(\frac{R(z_j, t_j)}{\xi}\right)^2
\end{equation}

\subsection{Anchor Loss ($\mathcal{L}_{\text{anchor}}$)}

MSE between PINN-predicted and HYDRUS-generated normalised pressure
heads at all 24,000 anchor points:

\begin{equation}
\mathcal{L}_{\text{anchor}} = \frac{1}{N_a}
\sum_{i=1}^{N_a}
\left(\tilde{\psi}_{\text{PINN}}(z_i, t_i) -
      \tilde{\psi}_{\text{HYDRUS}}(z_i, t_i)\right)^2
\end{equation}

\subsection{Initial Condition Loss ($\mathcal{L}_{\text{initial}}$)}

Enforces the pre-monsoon pressure head profile at $t = 0$ across all
1,000 depth nodes:

\begin{equation}
\mathcal{L}_{\text{initial}} = \frac{1}{N_i}
\sum_{m=1}^{N_i}
\left(\tilde{\psi}_{\text{PINN}}(z_m, 0) -
      \tilde{\psi}_{\text{HYDRUS}}(z_m, 0)\right)^2
\end{equation}

Without this term the network is free to predict any initial pressure
state, causing the temporal evolution of the Richards equation to
diverge from physical reality.

\subsection{Boundary Loss ($\mathcal{L}_{\text{boundary}}$)}

Matches the predicted Darcy flux at the surface ($z = z_{\min}$) to
the prescribed infiltration intensity at each of the 24 timesteps:

\begin{equation}
q_{\text{pred}} = -K(\psi)\left(\frac{\partial \psi}{\partial z}
+ 1\right), \qquad
\mathcal{L}_{\text{boundary}} = \frac{1}{N_b}
\sum_{k=1}^{N_b}
\left(q_{\text{pred}}(0, t_k) + I(t_k)\right)^2
\end{equation}

The current implementation uses $I(t_k) = K_s = 1.5 \times 10^{-5}$
m/s as a constant placeholder. Substituting the actual PERSIANN-CDR
daily rainfall record is identified as the immediate next step.

\subsection{Event-Based Failure Loss ($\mathcal{L}_{\text{failure}}$)}
\label{sec:fs}

A soft hinge penalty is applied at all 23,448 HYDRUS records where
FS $\leq 1.0$, encouraging the PINN to reproduce the geomechanical
criticality of the slope at depth. The Factor of Safety is computed
from the infinite-slope model:

\begin{equation}
FS(z,t) =
\frac{c' + \left(\gamma_{\text{eff}} z \cos^2\!\beta -
      \max(\psi,0)\,\gamma_w\right)\tan\phi'}
{\gamma_{\text{eff}} z \sin\beta\cos\beta}
\end{equation}

where $\gamma_{\text{eff}} = \gamma + \theta(z,t)\,\gamma_w$ is the
moisture-dependent dynamic unit weight. The penalty is:

\begin{equation}
\mathcal{L}_{\text{failure}} = \frac{1}{N_f}
\sum_{j=1}^{N_f}
\max\!\left(0,\; 1 - FS(z_j, t_j)\right)^2
\end{equation}

\section{Training Workflow}

\subsection{Two-Phase Optimisation Strategy}

\textbf{Phase I — Adam with StepLR (10,000 epochs):}
The Adam optimiser is used with an initial learning rate of
$\eta = 5 \times 10^{-4}$. A StepLR scheduler halves the learning
rate every 2,000 epochs, producing the schedule:
$5\times10^{-4} \to 2.5\times10^{-4} \to 1.25\times10^{-4}
\to 6.25\times10^{-5} \to 3.125\times10^{-5}$.
This decay prevents the oscillation of anchor and initial condition
losses that occurs at a fixed learning rate once the physics residual
has been minimised.

\textbf{Phase II — L-BFGS (50 outer steps $\times$ 100 iterations):}
L-BFGS with history size 50 and strong Wolfe line search refines
the solution using quasi-Newton curvature information. The optimiser
is called in 50 outer loops of 100 iterations each (5,000 total),
with early stopping if total loss improvement falls below $10^{-6}$
for five consecutive outer steps. This structure provides 50 visible
checkpoints and avoids the silent convergence failure that occurs
when all iterations run inside a single \texttt{.step()} call.

\subsection{Software and Hardware Environment}

\begin{itemize}
    \item \textbf{Framework:} PyTorch 2.0+, Python 3.12
    \item \textbf{Automatic differentiation:} PyTorch
    \texttt{autograd} computes $\partial\psi/\partial z$,
    $\partial\psi/\partial t$, and $\partial^2\psi/\partial z^2$
    at collocation points
    \item \textbf{Collocation sampling:} \texttt{scipy.stats.qmc.LatinHypercube}
    \item \textbf{Hardware:} CPU (Intel); no GPU required for
    inference after training
    \item \textbf{Normalisation:} All coordinates and outputs
    normalised to $[0, 1]$; physical de-normalisation applied
    inside the loss function before PDE evaluation
\end{itemize}
\chapter{Experimental Setup}

This project does not involve physical laboratory experiments or
field instrumentation. The ``experiment'' is a computational one:
a fully reproducible software pipeline that generates synthetic
training data, trains a Physics-Informed Neural Network, and
evaluates slope stability from the trained model. This chapter
describes the tools, configuration, and execution procedure that
constitute the experimental setup.

\section{Software and Hardware Environment}

\begin{table}[htbp]
\centering
\begin{tabular}{ll}
\hline
\textbf{Component} & \textbf{Specification} \\
\hline
Operating System    & Ubuntu 24.04 / macOS \\
Language            & Python 3.12 \\
Deep Learning       & PyTorch 2.0+ \\
Numerical Solver    & HYDRUS-1D v4.xx \\
Sampling            & SciPy 1.11+ (\texttt{stats.qmc.LatinHypercube}) \\
Data Processing     & NumPy, Pandas \\
Visualisation       & Matplotlib, Plotly \\
Dashboard (legacy)  & Streamlit \\
Dashboard (primary) & React 19 + FastAPI + Plotly.js \\
Hardware            & CPU (Intel/AMD); no GPU required \\
\hline
\end{tabular}
\caption{Software and hardware environment.}
\label{tab:env}
\end{table}

\section{Pipeline Stages}

The experimental pipeline runs in three sequential stages, each
corresponding to a script in the project codebase.

\subsection{Stage 1 — Data Ingestion}

The raw HYDRUS-1D output file (\texttt{Nod\_Inf.out}) is parsed
line by line. Each record containing a timestep header and
depth--pressure head pair is extracted and assembled into a
structured CSV (\texttt{artifacts/dataset/final\_data.csv})
with columns \texttt{[Time\_days, Depth\_m, Pressure\_Head, FS]}.
The Factor of Safety is computed at every point using the
infinite-slope model with the geotechnical parameters from
\texttt{params.yaml}. This stage produces 24,000 records across
24 daily timesteps and 1,000 depth nodes.

\subsection{Stage 2 — Data Loading and Batch Construction}

The processed CSV is converted into a serialised PyTorch batch
dictionary (\texttt{artifacts/dataset/data\_batch.pt}) containing
all tensors required for training:

\begin{itemize}
    \item \textbf{Anchor tensors} (\texttt{z\_data},
    \texttt{t\_data}, \texttt{psi\_data}): all 24,000 normalised
    HYDRUS records.

    \item \textbf{Collocation tensors} (\texttt{z\_coll},
    \texttt{t\_coll}): 10,000 points sampled via Latin Hypercube
    Sampling with \texttt{seed=42}, stored with
    \texttt{requires\_grad=True} for automatic differentiation.

    \item \textbf{Initial condition tensors} (\texttt{z\_ic},
    \texttt{t\_ic}, \texttt{psi\_ic}): 1,000 depth nodes at
    $t = t_{\min}$, enforcing the pre-monsoon pressure profile.

    \item \textbf{Boundary tensors} (\texttt{z\_bc},
    \texttt{t\_bc}, \texttt{target\_flux}): 24 surface points,
    one per timestep, with flux $= -K_s$.

    \item \textbf{Failure tensors} (\texttt{z\_fail},
    \texttt{t\_fail}): 23,448 points where FS $\leq 1.0$ and
    depth $> 0.1$ m.

    \item \textbf{Normalisation parameters}
    (\texttt{norm\_params}): \{$z_{\max}=55.0$ m,
    $t_{\max}=123.0$ days, $\psi_{\min}=-500.0$ m,
    $\psi_{\max}=-46.91$ m\}, saved alongside the batch and also
    written to \texttt{artifacts/model/norm\_params.json} after
    training.
\end{itemize}

\subsection{Stage 3 — Model Training}

The batch is loaded once into memory. The PINN (7 hidden layers,
64 neurons, tanh, Xavier initialisation, 29,377 parameters) is
instantiated and trained using the two-phase strategy described in
Section 5.5. Loss values for all five components are logged every
500 epochs to \texttt{logs/running\_logs.log} and serialised to
\texttt{artifacts/model/loss\_history.json} after completion. The
final model weights are saved to
\texttt{artifacts/model/pinn\_model.pt}.

\section{Configuration Files}

All experimental parameters are controlled through two YAML
configuration files, ensuring full reproducibility.

\textbf{\texttt{config/config.yaml}} defines artifact paths:
\begin{itemize}
    \item Raw data: \texttt{dataset/Nod\_Inf.out}
    \item Processed data: \texttt{artifacts/dataset/final\_data.csv}
    \item Training batch: \texttt{artifacts/dataset/data\_batch.pt}
    \item Trained model: \texttt{artifacts/model/pinn\_model.pt}
\end{itemize}

\textbf{\texttt{params.yaml}} defines all physical and
training parameters:
\begin{itemize}
    \item Architecture: \texttt{n\_hidden\_layers: 7},
    \texttt{neurons\_per\_layer: 64}, \texttt{activation: tanh}
    \item Training: \texttt{adam\_epochs: 10000},
    \texttt{lbfgs\_epochs: 5000},
    \texttt{learning\_rate: 0.0005}
    \item Loss weights: $\lambda_{\text{physics}} = 10^6$,
    $\lambda_{\text{anchor}} = 10$,
    $\lambda_{\text{initial}} = 10$,
    $\lambda_{\text{boundary}} = 1$,
    $\lambda_{\text{failure}} = 20$
    \item Geotechnical parameters: as listed in
    Table~\ref{tab:model_parameters}
\end{itemize}

\section{Reproducibility}

The entire pipeline is executed with a single command:

\begin{verbatim}
python main.py
\end{verbatim}

Latin Hypercube collocation points are generated with a fixed
\texttt{seed=42}, making the training batch identical across runs.
All artifact paths are configurable via \texttt{config.yaml}.
The trained model and normalisation parameters are versioned
together in \texttt{artifacts/model/} so that any prediction run
is guaranteed to use the correct de-normalisation boundaries.
\chapter{System Design}

This chapter presents the architectural and behavioural design of the
hybrid PINN framework through a series of diagrams covering the overall
system architecture, data flow, network structure, loss function
composition, and deployment interface.

\section{Overall System Architecture}

Figure~\ref{fig:system_arch} illustrates the end-to-end pipeline from
raw rainfall data to real-time slope stability output. The system is
organised into three subsystems: the numerical baseline (HYDRUS-1D
teacher model), the PINN training engine, and the inference and
deployment layer.

\begin{figure}[htbp]
    \centering
    \includegraphics[width=\linewidth]{Images/pinn_landslide_architecture.pdf}
    \caption{Overall system architecture of the hybrid PINN framework
    for rainfall-induced landslide prediction.}
    \label{fig:system_arch}
\end{figure}

\section{Data Flow Diagram}

Figure~\ref{fig:dataflow} shows the three-stage data pipeline:
Stage 1 parses the HYDRUS-1D output into a structured CSV; Stage 2
constructs the normalised PyTorch training batch containing five
distinct tensor sets; Stage 3 trains the PINN and serialises the
model weights and normalisation parameters.

\begin{figure}[htbp]
    \centering
    \includegraphics[width=\linewidth]{Images/dataflow.png}
    \caption{Three-stage data pipeline from HYDRUS-1D output to
    trained model artefacts.}
    \label{fig:dataflow}
\end{figure}

\section{PINN Network Architecture}

Figure~\ref{fig:pinn_arch} shows the feed-forward network structure.
The network accepts two normalised inputs — depth $\tilde{z}$ and
time $\tilde{t}$ — passes them through 7 fully connected hidden layers
of 64 neurons each with tanh activation, and produces a single
normalised pressure head output $\tilde{\psi}$. The same forward pass
is differentiated automatically to compute
$\partial\psi/\partial z$, $\partial\psi/\partial t$, and
$\partial^2\psi/\partial z^2$ for the Richards equation residual.

\begin{figure}[htbp]
    \centering
    \includegraphics[width=0.85\linewidth]{Images/logo.png}
    \caption{Physics-Informed Neural Network architecture:
    2-input, 7 $\times$ 64-neuron hidden layers (tanh),
    1-output feed-forward network with automatic differentiation
    for PDE enforcement.}
    \label{fig:pinn_arch}
\end{figure}

\section{Multi-Component Loss Function}

Figure~\ref{fig:loss_diagram} illustrates how the five loss
components are computed from their respective input sets and combined
into the total weighted loss that drives backpropagation.

\begin{figure}[htbp]
    \centering
    \includegraphics[width=\linewidth]{Images/loss_diagram.png}
    \caption{Multi-component loss function composition showing the
    five input sets, their corresponding loss terms, adopted weights,
    and weighted summation into $\mathcal{L}_{\text{total}}$.}
    \label{fig:loss_diagram}
\end{figure}

\section{Two-Phase Training Workflow}

Figure~\ref{fig:training_flow} shows the sequential optimisation
strategy. Phase I uses Adam with StepLR decay over 10,000 epochs for
global search. Phase II uses L-BFGS in 50 outer loops of 100
iterations each for quasi-Newton refinement, with early stopping when
loss improvement falls below $10^{-6}$ for five consecutive steps.

\begin{figure}[htbp]
    \centering
    \includegraphics[width=0.9\linewidth]{Images/training_flow.png}
    \caption{Two-phase training workflow: Adam with StepLR
    (Phase I) followed by L-BFGS with outer loop and early
    stopping (Phase II).}
    \label{fig:training_flow}
\end{figure}

\section{Inference Pipeline}

Figure~\ref{fig:inference} illustrates the post-training inference
flow. Given arbitrary depth and time query points alongside hourly
rainfall input, the trained PINN predicts $\psi(z,t)$, from which
Van Genuchten hydraulic properties and the infinite-slope Factor of
Safety are computed in under 50 milliseconds.

\begin{figure}[htbp]
    \centering
    \includegraphics[width=0.9\linewidth]{Images/inference.png}
    \caption{Inference pipeline from query inputs to Factor of
    Safety output via the trained PINN virtual sensor.}
    \label{fig:inference}
\end{figure}

\section{Dashboard Architecture}

Figure~\ref{fig:dashboard_arch} shows the deployment architecture
of the production dashboard. A FastAPI backend exposes 16 REST
endpoints that serve PINN inference, HYDRUS comparison, uncertainty
quantification, and sensitivity analysis. A React 19 + Vite frontend
consumes these endpoints across 17 interactive pages. A legacy
Streamlit dashboard shares the same inference engine via a common
\texttt{model\_inference.py} module.

\begin{figure}[htbp]
    \centering
    \includegraphics[width=\linewidth]{Images/pinn_landslide_architecture.svg}
    \caption{Deployment architecture: FastAPI backend with 16
    endpoints, React + Vite frontend with 17 pages, and legacy
    Streamlit interface sharing a common inference engine.}
    \label{fig:dashboard_arch}
\end{figure}

\subsection{Purpose of the UI Component}

The dashboard was developed to make the trained PINN usable as an
interactive engineering tool rather than only a command-line research
prototype. Its purpose is to allow a user to inspect pressure head and
Factor of Safety outputs, compare PINN results with HYDRUS-1D
reference data, explore parameter sensitivity, and export derived
results for reporting. This interface is important in the context of
the minor project because it demonstrates that the trained model can be
used for practical interpretation and not only for offline numerical
evaluation.

\subsection{Frontend Component Design}

The production frontend is implemented in the \texttt{ui-react}
module using React 19, Vite, React Router, and Plotly.js. The
application is organised into reusable components so that the same
layout and interaction patterns are used across all visual pages.

\begin{itemize}
    \item \textbf{Application shell:} The main application defines 17
    routes, each corresponding to a specific analysis page such as
    pressure head, Factor of Safety, HYDRUS comparison, uncertainty,
    and validation.

    \item \textbf{Shared layout:} A reusable sidebar groups pages into
    Core, Advanced, and Research sections. The sidebar is collapsible
    on desktop and opens as a mobile overlay on smaller screens.

    \item \textbf{Reusable UI elements:} Cards, page headers, metric
    tiles, tabs, spinners, and slider controls are implemented as
    shared components to keep the interface consistent.

    \item \textbf{Global state:} A shared context loads default
    geotechnical and normalisation parameters from the backend and
    propagates them to all pages. This allows parameter changes to be
    reflected immediately across the dashboard.

    \item \textbf{API abstraction:} All frontend requests are routed
    through a single API utility, which simplifies communication with
    the FastAPI service and keeps the page code focused on
    visualisation logic.
\end{itemize}

This modular structure improves maintainability and supports future
expansion of the dashboard without major changes to the overall
application design.

\subsection{Explanation of the Visualizer Part}

The visualizer is the most important part of the UI because it converts
PINN predictions into interpretable graphics for engineering analysis.
Rather than using a single chart, the system provides several
specialised visual modules, each intended to answer a different
question about slope behaviour.

\subsubsection*{Pressure Head Visualizer}

The pressure head page visualises the predicted hydraulic field
$\psi(z,t)$ in two complementary ways. First, it plots pressure head
against depth for selected timesteps so that the change in suction
profile through the monsoon period can be examined directly. Second,
it plots pressure head against time at a user-selected depth, making it
possible to observe temporal hydraulic response at a fixed location in
the slope. This page acts as the primary hydraulic interpretation
interface before stability is examined.

\subsubsection*{Factor of Safety Visualizer}

The Factor of Safety page is the main stability visualizer of the
system. It computes a full depth-time grid of FS values and displays:

\begin{itemize}
    \item a two-dimensional heatmap showing stable, marginal, and
    unstable regions over the full domain,
    \item summary statistics for safe, marginal, and unstable cells,
    and
    \item time-history curves of FS at user-selected depths.
\end{itemize}

This visualizer makes it possible to identify when instability starts,
which depth zone is most failure-prone, and whether safety improves or
degrades with time.

\subsubsection*{Time-Lapse Visualizer}

The animation page extends the visualizer by presenting the hydraulic
and stability response as a time-lapse sequence. Frames for all
available timesteps are pre-fetched from the backend, and the user can
play, pause, reset, scrub through the timeline, and switch between
pressure head and Factor of Safety modes. This feature is especially
useful for demonstrating temporal evolution during presentations and
for showing how hydraulic changes precede mechanical instability.

\subsubsection*{Validation and Comparison Visualizer}

The HYDRUS comparison page provides the validation-oriented part of the
visualizer. It includes overlay plots of HYDRUS and PINN profiles at
selected timesteps, depth-wise absolute error profiles, scatter plots
of predicted versus reference values, and tabulated statistics such as
$R^2$, RMSE, and MAE. This ensures that the dashboard does not only
display predictions, but also provides direct visual evidence of model
accuracy.

\subsubsection*{Advanced Visual Analytics}

Additional pages extend the visualizer beyond basic prediction plots.
These include PDE residual heatmaps for physics verification, Monte
Carlo uncertainty plots for stochastic FS analysis, critical slip depth
tracking, parameter sensitivity views, and export pages for structured
data download. Together, these modules turn the dashboard into an
exploratory decision-support tool rather than a passive display layer.

\subsection{Backend Support and Interaction Workflow}

The FastAPI backend supports the visualizer through dedicated endpoints
such as \texttt{/api/predict}, \texttt{/api/fs-grid},
\texttt{/api/hydrus-comparison}, \texttt{/api/pde-residual},
\texttt{/api/uncertainty}, and \texttt{/api/critical-slip}. These
endpoints receive user-selected inputs from the frontend, evaluate the
trained PINN or the required post-processing routine, and return
structured JSON arrays that are rendered directly as Plotly charts,
metric cards, and statistics tables.

The interaction workflow is therefore:

\begin{enumerate}
    \item the dashboard loads default parameters and available
    timesteps from the backend,
    \item the user selects a visual analysis page and adjusts the
    required inputs,
    \item the frontend sends the corresponding API request,
    \item the backend evaluates the PINN model or the requested
    analysis routine, and
    \item the response is rendered as an interactive chart, table, or
    summary card.
\end{enumerate}

This workflow allows rapid experimentation and makes the project easier
to demonstrate during evaluation. The visualizer is therefore not an
optional interface layer, but an essential part of the project because
it improves interpretability, validation, and practical usability.

\section{Factor of Safety Computation Model}

Figure~\ref{fig:fs_model} summarises the infinite-slope stability
model used to convert PINN-predicted pressure head into Factor of
Safety at every queried $(z, t)$ point, showing the inputs, the
dynamic unit weight calculation, and the FS formula.

\begin{figure}[htbp]
    \centering
    \includegraphics[width=0.85\linewidth]{Images/fs_model.png}
    \caption{Infinite-slope Factor of Safety computation model
    driven by PINN-predicted $\psi(z,t)$ and static geotechnical
    parameters.}
    \label{fig:fs_model}
\end{figure}
\chapter{Results and Discussion}

\section{Results}

\subsection{HYDRUS-1D Numerical Baseline}

The HYDRUS-1D simulation of the Jure slope under the 2014 monsoon
produced the pseudo-data that anchors the entire PINN training.
Table~\ref{tab:hydrus_summary} summarises the dataset characteristics.

\begin{table}[htbp]
\centering
\begin{tabular}{ll}
\hline
\textbf{Property} & \textbf{Value} \\
\hline
Total pseudo-data points      & 24,000 \\
Temporal coverage             & 24 timesteps, 0 -- 123 days \\
Spatial coverage              & 1,000 depth nodes, 0 -- 55 m \\
Pressure head range           & $-500.0$ to $-46.9$ m \\
Factor of Safety range        & 0.600 to 9.73 \\
Failure points (FS $\leq$ 1.0, $z > 0.1$ m) & 23,448 (97.7\%) \\
\hline
\end{tabular}
\caption{HYDRUS-1D pseudo-data summary.}
\label{tab:hydrus_summary}
\end{table}

The pressure head remained entirely negative throughout all 123
simulated days, indicating that the slope stayed in an unsaturated
state with no positive pore water pressures developing.
Table~\ref{tab:psi_evolution} shows the progressive reduction in
mean suction across the simulation period, which captures the
antecedent moisture buildup that preceded the August 2 failure.

\begin{table}[htbp]
\centering
\begin{tabular}{rrrr}
\hline
\textbf{Time (days)} & \textbf{$\psi_{\min}$ (m)} &
\textbf{$\psi_{\text{mean}}$ (m)} & \textbf{$\psi_{\max}$ (m)} \\
\hline
0   & $-500.00$ & $-275.06$ & $-50.00$ \\
1   & $-196.93$ & $-134.38$ & $-51.73$ \\
10  & $-139.16$ & $-82.44$  & $-55.00$ \\
30  & $-131.01$ & $-80.88$  & $-56.33$ \\
60  & $-123.63$ & $-78.56$  & $-54.30$ \\
96  & $-116.97$ & $-73.77$  & $-49.03$ \\
123 & $-112.69$ & $-71.61$  & $-46.91$ \\
\hline
\end{tabular}
\caption{Pressure head evolution across the 55 m column at
selected timesteps.}
\label{tab:psi_evolution}
\end{table}

The sharpest change occurs between day 0 and day 1, where the column
mean suction reduces by 141 m — from $-275$ m to $-134$ m —
reflecting the rapid infiltration response at monsoon onset.
From day 10 onward the profile evolves gradually, with mean suction
declining from $-82$ m to $-72$ m over the remaining 113 days.
This slow progressive saturation represents the antecedent loading
that primed the slope for failure.

\subsection{PINN Training Convergence}

Table~\ref{tab:training_loss} reports the per-component loss values
at selected training checkpoints. Figure~\ref{fig:loss_curve} presents
the corresponding loss curves for the full 10,000-epoch Adam phase.

\begin{table}[htbp]
\centering
\begin{tabular}{rrrrrr}
\hline
\textbf{Epoch} & \textbf{Total} & \textbf{Physics (raw)} &
\textbf{Anchor} & \textbf{Initial} & \textbf{Failure} \\
\hline
1       & $1.390 \times 10^{2}$ & $1.259 \times 10^{-4}$ & 0.7122 & 0.3129 & 0.1432 \\
500     & $4.952$               & $1.321 \times 10^{-7}$ & 0.1387 & 0.0567 & 0.1433 \\
1,000   & $4.352$               & $3.940 \times 10^{-8}$ & 0.0590 & 0.0855 & 0.1434 \\
5,000   & $4.233$               & $2.678 \times 10^{-9}$ & 0.0708 & 0.0655 & 0.1434 \\
10,000  & $4.231$               & $1.775 \times 10^{-8}$ & 0.0713 & 0.0633 & 0.1434 \\
\hline
\end{tabular}
\caption{Per-component training loss at selected Adam epochs.}
\label{tab:training_loss}
\end{table}

Three convergence phases are evident. During epochs 1--500 the total
loss fell 96.4\% from 139.0 to 4.95, driven by the physics term which
reduced by \textbf{99.9\%} as the Richards equation residual was
rapidly eliminated. During epochs 500--1,000 the anchor and initial
condition losses became the dominant signals, with anchor dropping
from 0.139 to 0.059 and IC from 0.057 to 0.086. From epoch 1,000
onward the total loss plateaued between 4.22 and 4.37, oscillating
without further convergence. Table~\ref{tab:convergence_summary}
summarises the total reduction achieved by each component.

\begin{table}[htbp]
\centering
\begin{tabular}{lrr}
\hline
\textbf{Loss component} & \textbf{Epoch 1 value} &
\textbf{Reduction (\%)} \\
\hline
Physics (Richards residual) & $1.259 \times 10^{-4}$ & \textbf{99.99} \\
Anchor (HYDRUS data fit)    & 0.7122                  & 89.98 \\
Initial condition           & 0.3129                  & 79.78 \\
Failure (hinge penalty)     & 0.1432                  & 0.14 \\
\hline
\end{tabular}
\caption{Loss component reductions from epoch 1 to epoch 10,000.}
\label{tab:convergence_summary}
\end{table}

The failure loss remained unchanged at 0.1434 throughout all training
phases. This stability is a physically meaningful result: it
demonstrates that the slope at Jure is in a state of persistent
geomechanical criticality below 1.2 m depth, a condition that
reflects real physical parameters and cannot be altered by
optimisation.

\subsection{Pressure Head Prediction Accuracy}

The trained PINN predicted pressure heads ranging from $-500.2$ m
to $-85.3$ m across the 72-hour inference scenario, entirely within
the training domain of $-500.0$ to $-46.9$ m. All predicted values
remained negative, confirming the absence of spurious positive pore
pressures.

Table~\ref{tab:accuracy} reports error metrics from point-by-point
comparison of PINN predictions against HYDRUS reference values at
$t = 0$ across the 0--3 m depth range.

\begin{table}[htbp]
\centering
\begin{tabular}{ll}
\hline
\textbf{Metric} & \textbf{Value} \\
\hline
Root Mean Square Error (RMSE)              & 0.724 m \\
Mean Absolute Error (MAE)                  & 0.585 m \\
Maximum absolute error                     & 1.526 m \\
Relative RMSE (\% of 453 m range)          & \textbf{0.16\%} \\
Predictions within 1.0 m of HYDRUS         & 78\% (31 of 40) \\
Predictions within 2.0 m of HYDRUS         & 100\% (40 of 40) \\
\hline
\end{tabular}
\caption{PINN prediction accuracy vs HYDRUS-1D reference at $t = 0$.}
\label{tab:accuracy}
\end{table}

A relative RMSE of 0.16\% over a 453 m pressure head range confirms
that the network has learned the spatial structure of the field
accurately. All 40 prediction points are within 2 m of the HYDRUS
reference, and 78\% are within 1 m. The largest errors are
concentrated at the surface zone ($z < 0.5$ m), where the pressure
head undergoes the sharpest transient changes under rainfall forcing.
At depths beyond 1.5 m, absolute errors fall consistently below
0.5 m. Figure~\ref{fig:psi_comparison} shows the overlay of PINN
predictions against HYDRUS profiles at selected timesteps.

\subsection{Factor of Safety Analysis}

\subsubsection{Critical Stability Transition Depth}

A primary quantitative result is the identification of the critical
stability transition depth — the depth at which Factor of Safety
crosses 1.0. In the HYDRUS training data at $t = 0$, FS crosses
unity between $z = 1.211$ m (FS $= 1.006$) and $z = 1.266$ m
(FS $= 0.988$). The PINN predictions locate the same crossover
between $z = 1.162$ m (FS $= 1.005$) and $z = 1.186$ m
(FS $= 0.997$) — a spatial discrepancy of \textbf{0.05 m (4.1\%)}
from the HYDRUS reference. This close agreement, produced
independently by the PINN, confirms that the network correctly
learned both the pressure head field and its geomechanical
implications.

\subsubsection{Depth-Zone Stability Distribution}

Table~\ref{tab:fs_zones} summarises the Factor of Safety statistics
across the 2,880 prediction points of the 72-hour inference scenario,
organised by depth zone. Figure~\ref{fig:fs_heatmap} presents the
corresponding FS heatmap over the full (depth, time) domain.

\begin{table}[htbp]
\centering
\begin{tabular}{lrrrrc}
\hline
\textbf{Depth zone (m)} & \textbf{$n$} & \textbf{FS$_{\min}$} &
\textbf{FS$_{\text{mean}}$} & \textbf{FS$_{\max}$} &
\textbf{Stable (\%)} \\
\hline
0.0 -- 0.5 & 360 & 1.827 & 2.527 & 3.611 & 100.0 \\
0.5 -- 1.0 & 864 & 1.073 & 1.284 & 1.546 & 100.0 \\
1.0 -- 1.5 & 360 & 0.924 & 0.959 & 1.005 & 6.1   \\
1.5 -- 2.0 & 504 & 0.844 & 0.860 & 0.894 & 0.0   \\
2.0 -- 3.0 & 792 & 0.753 & 0.781 & 0.822 & 0.0   \\
\hline
\end{tabular}
\caption{PINN-predicted Factor of Safety statistics by depth zone
across all 72 simulated hours.}
\label{tab:fs_zones}
\end{table}

The near-surface zone (0--1.0 m) maintains persistent stability with
FS between 1.07 and 3.61 throughout all 72 simulated hours.
The transition zone (1.0--1.5 m) straddles the critical threshold,
with only 6.1\% of predictions remaining stable. All zones below
1.5 m show FS $< 1.0$ at every prediction point, with a minimum
of 0.753 at the deepest sampled depth.

\subsection{Inference Performance}

The trained PINN evaluates $\psi(z,t)$ and FS at 2,880 query points
(40 depths $\times$ 72 hours) in under 50 milliseconds on a standard
CPU. Table~\ref{tab:performance} compares this against conventional
approaches.

\begin{table}[htbp]
\centering
\begin{tabular}{lllc}
\hline
\textbf{Method} & \textbf{Evaluation time} &
\textbf{Points} & \textbf{Real-time capable} \\
\hline
PINN (this work)   & $<$ 50 ms     & 2,880 arbitrary & Yes \\
HYDRUS-1D          & 30 -- 120 s   & 1,000 fixed nodes & No \\
FEM (PLAXIS)       & 5 -- 30 min   & $\sim$10,000 elements & No \\
\hline
\end{tabular}
\caption{Computational performance comparison.}
\label{tab:performance}
\end{table}

\section{Discussion}

\subsection{Physical Interpretation of the Stability Results}

The consistent FS $< 1.0$ below 1.2 m depth throughout all
timesteps is a correct geomechanical result rather than a modelling
artefact. The slope angle $\beta = 42°$ substantially exceeds the
effective friction angle $\phi' = 28°$. For a purely frictional
soil, the infinite-slope FS reduces to $\tan\phi'/\tan\beta =
\tan(28°)/\tan(42°) = 0.598$, which matches exactly the minimum
FS observed in the HYDRUS dataset (0.5997). This identity confirms
that the model is computing correct geomechanics: at depths beyond
approximately 1.2 m, cohesion $c' = 5.0$ kPa is insufficient to
offset the gravitational shear stress, and the slope exists in a
state of persistent geomechanical criticality.

This finding is consistent with published back-analyses of the Jure
site (Bray et al., 2023; Tien et al., 2021), which characterise it
as a slope in critical equilibrium sustained by matric suction.
Progressive suction loss under sustained rainfall reduced the apparent
cohesion contribution and triggered collapse — a mechanism the PINN
correctly reproduces by identifying the suction-dependent stability
boundary at 1.2 m.

\subsection{Significance of the 99.99\% Physics Residual Reduction}

The physics loss reduction from $1.259 \times 10^{-4}$ to
$1.775 \times 10^{-8}$ — a 99.99\% decrease — demonstrates that
the network has learned to satisfy the Richards equation to near
machine precision. This is the core claim of any Physics-Informed
Neural Network: that physical constraints are not just loosely
encouraged but actively enforced in the solution. The normalised
residual scale of $1.775 \times 10^{-8}$ represents $1.34 \times
10^{-4}$ of the characteristic Richards residual magnitude for this
soil, confirming near-perfect PDE satisfaction across all 10,000
collocation points.

\subsection{Loss Plateau Analysis}

The Adam optimiser plateau from epoch 1,000 onward, with total loss
oscillating between 4.22 and 4.37, is a structural feature of the
loss landscape rather than a training failure. The plateau floor is
determined by the failure loss, which contributes $\lambda_{\text{failure}}
\times 0.1434 = 20 \times 0.1434 = 2.868$ to the total — representing
the irreducible geomechanical criticality of the Jure slope. The
remaining $\approx 1.37$ of total loss is reducible through
improved training (StepLR decay and L-BFGS refinement in the next
run), and is expected to push anchor RMSE from the current 121 m
(26.7\% of range) below 50 m.

\subsection{Virtual Sensor Validation}

The sub-50 ms inference time, combined with the 0.16\% relative
RMSE and 4.1\% transition depth accuracy, validates the virtual
sensor concept. The PINN delivers physically consistent
pressure head and stability estimates at any (depth, time) query
three to four orders of magnitude faster than the HYDRUS solver
it was trained on. This is the defining practical advantage of
the framework: once trained, it removes the computational barrier
to real-time slope monitoring entirely.

\section{Testing, Evaluation, and Validation}

\subsection{Physics Consistency Tests}

Three physics consistency checks were applied to the trained model:

\begin{enumerate}
    \item \textbf{Sign check:} All 2,880 predicted pressure head
    values are negative ($-500.2$ to $-85.3$ m), confirming the
    model correctly represents an unsaturated slope with no
    spurious positive pore pressures.

    \item \textbf{Monotonicity check:} Predicted pressure head
    profiles at each timestep decrease monotonically from the
    surface toward depth (suction reduces with depth), consistent
    with downward infiltration physics. No inversions were observed.

    \item \textbf{PDE residual check:} The normalised Richards
    equation residual at convergence is $1.775 \times 10^{-8}$
    (99.99\% reduction from initialisation), confirming near-perfect
    PDE satisfaction across all collocation points.
\end{enumerate}

\subsection{Quantitative Accuracy vs HYDRUS-1D}

The primary quantitative validation compares PINN predictions against
the HYDRUS-1D reference profiles. At $t = 0$ across the 0--3 m
depth range: RMSE $= 0.724$ m, MAE $= 0.585$ m, relative error
$= 0.16\%$, and 100\% of predictions within 2 m of the reference.
The critical stability transition depth is reproduced to within
0.05 m (4.1\%) of the HYDRUS analytical value.

\subsection{Hydrological Performance Ratings}

Performance is assessed against the standard Moriasi et al.\ (2007)
hydrological model benchmark criteria, which classify model
performance as Very Good (NSE $> 0.75$), Good (NSE $> 0.65$),
Satisfactory (NSE $> 0.50$), or Unsatisfactory (NSE $\leq 0.50$).
These metrics are computed on the full training dataset through the
HYDRUS Comparison page of the dashboard. Table~\ref{tab:hydro_metrics}
presents the results.

\begin{table}[htbp]
\centering
\begin{tabular}{llll}
\hline
\textbf{Metric} & \textbf{Value} & \textbf{Rating} &
\textbf{Threshold} \\
\hline
Nash-Sutcliffe Efficiency (NSE) & [from dashboard] & --- & $> 0.75$ = Very Good \\
Kling-Gupta Efficiency (KGE)    & [from dashboard] & --- & $> 0.75$ = Very Good \\
$R^2$                           & [from dashboard] & --- & $> 0.75$ = Very Good \\
RMSE (overall)                  & [from dashboard] & --- & --- \\
PBIAS (\%)                      & [from dashboard] & --- & $< \pm 10$ = Very Good \\
\hline
\end{tabular}
\caption{Hydrological validation metrics from the dashboard
Validation page. Values to be recorded from
\texttt{/api/validation} endpoint.}
\label{tab:hydro_metrics}
\end{table}

\subsection{Objective Achievement Summary}

Table~\ref{tab:objectives} provides a final assessment of each
project objective against the results obtained.

\begin{table}[htbp]
\centering
\begin{tabular}{p{5.5cm}lp{5cm}}
\hline
\textbf{Objective} & \textbf{Status} & \textbf{Evidence} \\
\hline
PINN with Richards equation and Van Genuchten model
  & Achieved
  & Physics residual reduced 99.99\% \\
Hybrid training with HYDRUS anchor points
  & Achieved
  & 24,000 anchor pts; anchor loss reduced 90\% \\
Predict $\psi(z,t)$ and characterise stability
  & Achieved
  & RMSE $= 0.16\%$; transition depth within 4.1\% \\
Production virtual sensor and dashboard
  & Achieved
  & $<50$ ms inference; 17-page dual dashboard;
    16 REST endpoints \\
\hline
\end{tabular}
\caption{Project objectives vs achieved results.}
\label{tab:objectives}
\end{table}

\subsection{Figure References}

The following figures support the results presented in this chapter
and should be placed at the appropriate positions in the text:

\begin{itemize}
    \item \textbf{Figure~\ref{fig:loss_curve}:} Training loss curves
    (total, physics, anchor, IC, failure) over 10,000 Adam epochs —
    from \texttt{artifacts/model/loss\_history.json} via the
    Training Loss dashboard page.

    \item \textbf{Figure~\ref{fig:psi_comparison}:} PINN vs HYDRUS
    pressure head overlay at selected timesteps (days 0, 30, 60, 96,
    123) — from the HYDRUS Comparison dashboard page.

    \item \textbf{Figure~\ref{fig:fs_heatmap}:} Factor of Safety
    heatmap over depth $\times$ time (72-hour scenario) —
    from \texttt{artifacts/results/prediction\_plots.png}.

    \item \textbf{Figure~\ref{fig:psi_heatmap}:} Pressure head
    heatmap $\psi(z,t)$ over depth $\times$ time —
    from \texttt{artifacts/results/prediction\_plots.png}.
\end{itemize}
\chapter{Conclusions}

This project developed a Hybrid Physics-Informed Neural Network
framework for rainfall-induced landslide prediction in data-scarce
regions, demonstrated on the 2014 Jure landslide, Sindhupalchok,
Nepal. The framework solves the Richards equation for unsaturated
soil moisture dynamics and computes the Factor of Safety using the
infinite-slope stability model, without requiring any in-situ
instrumentation.

The PINN successfully learned the pressure head field $\psi(z,t)$
from 24,000 HYDRUS-1D pseudo-data points, achieving a 99.99\%
reduction in the Richards equation residual and a relative RMSE
of 0.16\% against the HYDRUS reference. The critical stability
transition depth was independently identified at $z = 1.16$--$1.19$
m, agreeing with the HYDRUS analytical value to within 4.1\%. The
trained model evaluates $\psi$, hydraulic properties, and Factor
of Safety at any queried depth and time in under 50 milliseconds —
three to four orders of magnitude faster than the HYDRUS solver it
was trained on — and is deployed through a 17-page interactive
dual dashboard with 16 REST API endpoints.

The results confirm that the Jure slope exists in a state of
persistent geomechanical criticality below 1.2 m depth, sustained
only by matric suction — a finding consistent with published
back-analyses of the site and providing a quantitative explanation
for its acute sensitivity to sustained monsoon rainfall. The
framework establishes the technical foundation for a low-cost
virtual sensor deployable at unmonitored Himalayan slopes where
conventional instrumentation is economically and logistically
infeasible.
\chapter{Limitations and Future Enhancements}

\section{Current Limitations}

\subsection{One-Dimensional Flow Simplification}

The framework models only 1D vertical infiltration, governed by the
Richards equation in the $z$-direction. The actual Jure landslide
involved lateral preferential flow through a heavily jointed
phyllite--schist formation with four primary joint sets
\cite{bray2023jure}. Slope-parallel drainage, lateral pressure
redistribution, and preferential pathways through fractures cannot
be represented in the current formulation. This simplification is
most consequential at depth ($z > 10$ m), where lateral flow
dominates the pre-failure pore pressure regime.

\subsection{Synthetic Training Data}

The anchor dataset comes entirely from HYDRUS-1D simulation rather
than in-situ sensor measurements. Model accuracy is therefore
bounded by the fidelity of the HYDRUS calibration. Uncertainty in
the geotechnical parameters — particularly $K_s$, $\alpha$, and
$n$ for the Kuncha Formation — propagates directly into the pseudo-
data and from there into the PINN. The current anchor RMSE of
121 m (26.7\% of the 453 m pressure head range) reflects this
inherited uncertainty.

\subsection{Surface Boundary Flux Placeholder}

The surface boundary condition currently uses a constant flux of
$-K_s = -1.5 \times 10^{-5}$ m/s at every timestep, rather than
the actual time-varying PERSIANN-CDR daily rainfall intensities.
This means the boundary loss term ($\lambda_{\text{boundary}} = 1$)
enforces an incorrect constraint and contributes negligibly to
training. The true rainfall-driven transient response at the surface
is learned implicitly through the anchor loss rather than explicitly
through the boundary condition.

\subsection{Adam Oscillation at Plateau}

The anchor and initial condition losses oscillated between 0.059 and
0.088 from epoch 1,000 to epoch 10,000 without converging, due to
the learning rate being too large relative to the flat loss landscape
once the physics residual was minimised. The planned StepLR
scheduler and L-BFGS outer loop fix have been implemented but not
yet applied to a full retrain.

\subsection{Fixed Failure Plane}

The event-based failure loss uses all HYDRUS records with
FS $\leq 1.0$ at depth $> 0.1$ m as failure points (23,448 points).
This broad criterion does not specifically target the documented
Jure failure plane at $z = 20$--$28$ m, and the geotechnical
parameters adopted place FS $< 1.0$ everywhere below 1.2 m —
meaning the failure loss does not distinguish between failure at
1.2 m and failure at 25 m. A targeted failure plane specification
would produce a more precise stability model.

\section{Future Enhancements}

\subsection{Real Rainfall Integration}

The most impactful immediate improvement is substituting the
constant $-K_s$ surface flux with the actual PERSIANN-CDR daily
rainfall record, converted to infiltration flux and capped at $K_s$
(Green-Ampt approximation). This requires adding a rainfall array
column to the processed CSV and mapping it to the 24 boundary
timesteps in the batch. Once activated, $\lambda_{\text{boundary}}$
can be raised to 5.0, making the surface response physically driven
by observed precipitation and enabling the model to reproduce
event-specific pore pressure dynamics.

\subsection{Targeted Failure Plane and Parameter Recalibration}

Recalibrating the geotechnical parameters to produce an initial
FS slightly above 1.0 at a shallower failure plane ($z = 1$--$3$ m)
would allow a clear rainfall-triggered FS decline from stable to
unstable to be demonstrated and timestamped. Increasing $c'$ to
8--10 kPa, or reducing $\beta$ to 35°--38°, would shift the
stability boundary below the full depth profile and allow the
failure timing objective to be achieved. This is the standard
approach in published PINN-based landslide studies
\cite{depina2022pinn, krkac2023pinn}.

\subsection{Inverse Problem — Hydraulic Parameter Discovery}

The trained PINN establishes the forward model infrastructure
required for inverse parameter identification. By treating $K_s$,
$\alpha$, and $n$ as learnable parameters and minimising the
physics loss against observed pore pressure or slope displacement
data, the framework can discover unknown soil hydraulic properties
from past failure records \cite{wang2025pinns, niu2025inverse}.
This transforms the system from a forward simulator into a
calibration tool, removing the current dependence on literature-
derived parameters for sites where no back-analysis exists.

\subsection{2D and 3D Extension}

Extending the governing PDE to the 2D Richards equation
($\partial\theta/\partial t = \nabla \cdot [K(\psi)(\nabla\psi +
\mathbf{e}_z)]$) would capture lateral flow and slope-parallel
drainage. The PINN architecture requires only an additional input
neuron (horizontal coordinate $x$) and additional collocation
dimensions — the loss function formulation is otherwise unchanged.
This extension is directly supported by the PyTorch autograd
framework already in use \cite{karniadakis2021pinns}.

\subsection{Transfer Learning to Other Himalayan Slopes}

The trained Jure model can serve as a pre-trained initialisation
for PINN training on adjacent or geologically similar slopes,
reducing the training data and computational cost required for
each new site \cite{giri2019transfer}. Transfer learning from
a well-calibrated source slope to a data-scarce target slope
would allow the framework to scale across the Himalayan region
without requiring a full HYDRUS calibration for every site —
directly addressing the deployment challenge for a regional
Early Warning System.

\subsection{Automated StepLR and Adaptive Loss Weighting}

Incorporating automatic loss weight adaptation — such as the
self-adaptive weighting scheme of Wang et al.\ (2022), which
dynamically adjusts $\lambda$ values based on gradient magnitudes
— would eliminate the manual weight tuning required in the current
implementation. Combined with the StepLR scheduler already
implemented, this would allow the framework to be deployed on
new sites without site-specific weight engineering.
\addcontentsline{toc}{section}{References}
\renewcommand{\bibname}{References}
\bibliographystyle{unsrt}
\bibliography{ref}
\chapter*{Appendices}
\addcontentsline{toc}{section}{Appendices}
It should contain the information that is somehow missing in the main body and it presents extra information about your project like program source code, raw data, and so on.
\end{document}
